% =====================================================================
% TEMPLATE — Rapport corporate Brantham Partners
% Style : LaTeX article, sobre, monochrome, serif Computer Modern
% Inspiration : Dossier_Complet.pdf (avr. 2026)
% Usage : remplir les champs <<...>>, compiler avec pdflatex
% =====================================================================

\documentclass[11pt,a4paper]{article}

\usepackage[utf8]{inputenc}
\usepackage[T1]{fontenc}
\usepackage[french]{babel}
\usepackage{lmodern}
\usepackage[a4paper,margin=2.5cm]{geometry}
\usepackage{enumitem}
\usepackage{booktabs}
\usepackage{array}
\usepackage{longtable}
\usepackage{tabularx}
\usepackage{hyperref}
\usepackage{titlesec}
\usepackage{fancyhdr}
\usepackage{setspace}

\hypersetup{
  colorlinks=false,
  pdfborder={0 0 0}
}

% Espacements titres : sobres, pas de couleurs
\titleformat{\section}{\Large\bfseries}{\thesection.}{0.6em}{}
\titleformat{\subsection}{\large\bfseries}{\thesubsection}{0.6em}{}
\titleformat{\subsubsection}{\normalsize\bfseries}{\thesubsubsection}{0.6em}{}

\titlespacing*{\section}{0pt}{2.2ex plus 1ex minus .2ex}{1.2ex plus .2ex}
\titlespacing*{\subsection}{0pt}{1.6ex plus 1ex minus .2ex}{0.8ex plus .2ex}

% Pied de page : numéro centré
\pagestyle{plain}

% Justification + interligne aéré
\setlength{\parskip}{0.5em}
\setlength{\parindent}{0pt}
\onehalfspacing

% =====================================================================
% MÉTADONNÉES — à remplir pour chaque dossier
% =====================================================================
\newcommand{\rapportTitre}{<<TITRE DU RAPPORT>>}
\newcommand{\rapportSousTitre}{<<Sous-titre / nature du rapport>>}
\newcommand{\rapportContexte}{Brantham Partners --- Conseil M\&A distressed}
\newcommand{\rapportDate}{<<AAAA-MM-JJ>>}
\newcommand{\rapportConfidentiel}{Confidentiel --- Usage interne et client}

\begin{document}

% =====================================================================
% PAGE DE GARDE — minimaliste, centrée verticalement
% =====================================================================
\thispagestyle{empty}
\vspace*{\fill}
\begin{center}
{\Large\bfseries \rapportTitre}\\[1em]
{\normalsize \rapportSousTitre}\\[3em]
{\small \rapportContexte}\\[1em]
{\small \rapportDate}\\[1em]
{\footnotesize \itshape \rapportConfidentiel}
\end{center}
\vspace*{\fill}
\newpage

% =====================================================================
% SOMMAIRE
% =====================================================================
\tableofcontents
\newpage

% =====================================================================
% AVANT-PROPOS (optionnel)
% =====================================================================
\section*{Avant-propos}
\addcontentsline{toc}{section}{Avant-propos}

<<Quelques lignes sur l'objet du rapport, son périmètre, ses limites, et la
posture analytique adoptée. Ne pas excéder 10 lignes.>>

\textbf{Utilisation recommandée.}
\begin{enumerate}[leftmargin=2em]
  \item <<Première lecture pour vue d'ensemble.>>
  \item <<Mémoriser les chiffres-clés et les drapeaux rouges.>>
  \item <<Se référer aux sources en annexe avant tout call client.>>
\end{enumerate}

\textbf{Structure du rapport.}
\begin{itemize}[leftmargin=2em]
  \item \textbf{Partie I} : <<...>>
  \item \textbf{Partie II} : <<...>>
  \item \textbf{Partie III} : <<...>>
\end{itemize}

\newpage

% =====================================================================
% CORPS DU RAPPORT
% =====================================================================

\section{Synthèse exécutive}

<<3 à 5 paragraphes de synthèse. Reprendre les 3 messages clés et la
recommandation finale en une seule phrase.>>

\textbf{Recommandation.} <<Une phrase, sans nuance superflue.>>

\subsection{Chiffres clés}

\begin{tabular}{ll}
\toprule
<<Indicateur 1>> & <<Valeur>> \\
<<Indicateur 2>> & <<Valeur>> \\
<<Indicateur 3>> & <<Valeur>> \\
\bottomrule
\end{tabular}

\section{Périmètre et structure du dossier}

\subsection{Entités concernées}

\textbf{<<Entité 1>>.} <<SIREN, forme juridique, capital, NAF, siège, date de
création. Une phrase de description du rôle dans le groupe.>>

\textbf{<<Entité 2>>.} <<Idem.>>

\subsection{Procédure collective}

\textbf{Tribunal.} <<Tribunal compétent, date du jugement d'ouverture, date de
cessation des paiements, période d'observation.>>

\textbf{Administrateur judiciaire.} <<Cabinet, nom du mandataire, adresse, contact.>>

\textbf{Mandataire judiciaire.} <<Idem.>>

\textbf{DLDO (date limite de dépôt des offres).} <<JJ/MM/AAAA --- horaire>>.

\section{Analyse financière}

\subsection{Compte de résultat consolidé}

\begin{tabular}{lrrrr}
\toprule
M\euro{} & 2022 & 2023 & 2024 & 2025 \\
\midrule
Chiffre d'affaires & & & & \\
EBITDA & & & & \\
Marge EBITDA & & & & \\
Résultat net & & & & \\
\bottomrule
\end{tabular}

\textbf{Lecture.} <<Diagnostic en 2 à 3 phrases. Souligner la rupture observée.>>

\subsection{Bilan et passif estimé}

\textbf{Dette financière.} <<Estimation de la dette senior, mezzanine, OBSA/OC.>>

\textbf{Passif fournisseurs.} <<Indicateurs Rubypayeur, retards observés.>>

\textbf{Passif social.} <<Effectif, masse salariale annuelle, exposition AGS.>>

\textbf{Passif fiscal et social (URSSAF).} <<Estimation arriérés.>>

\textbf{Passif total estimé.} <<Fourchette en M\euro{}.>>

\subsection{Tour de table et historique capitalistique}

<<Reconstituer le ou les tours de table successifs : date, sponsors, dette,
montants. Identifier la cap table actuelle (estimation si non publique).>>

\section{Valorisation et stratégie d'offre}

\subsection{Multiples sectoriels de référence}

\begin{itemize}[leftmargin=2em]
  \item <<Comparable 1 --- multiple appliqué>>
  \item <<Comparable 2 --- multiple appliqué>>
\end{itemize}

\subsection{Estimation de valeur d'entreprise}

\begin{tabular}{lrrr}
\toprule
Composante & Bear & Base & Bull \\
\midrule
<<Actif 1>> & & & \\
<<Actif 2>> & & & \\
<<Actif 3>> & & & \\
\midrule
\textbf{Total EV} & & & \\
\bottomrule
\end{tabular}

\subsection{Recommandation de prix d'offre}

\textbf{Règle interne Brantham.} Plafond \(=\) 12\,\% du CA de référence en
plan de cession.

\textbf{Offre cible recommandée.} <<Fourchette en k\euro{} ou M\euro{}, périmètre,
exclusions, conditions de reprise sociale.>>

\section{Drapeaux rouges}

\subsection{Criticité 1 --- bloquants}

\begin{enumerate}[leftmargin=2em]
  \item \textbf{<<Drapeau 1>>.} <<Description, impact, action correctrice.>>
  \item \textbf{<<Drapeau 2>>.} <<Idem.>>
\end{enumerate}

\subsection{Criticité 2 --- risques sérieux}

\begin{enumerate}[leftmargin=2em,start=3]
  \item \textbf{<<Drapeau 3>>.}
  \item \textbf{<<Drapeau 4>>.}
\end{enumerate}

\subsection{Criticité 3 --- à clarifier}

\begin{enumerate}[leftmargin=2em,start=5]
  \item \textbf{<<Drapeau 5>>.}
\end{enumerate}

\section{Cartographie repreneurs}

\subsection{Tier 1 --- outreach immédiat}

\begin{tabular}{p{4cm}p{3cm}p{6cm}}
\toprule
Cible & Dirigeant & Logique stratégique \\
\midrule
<<Repreneur 1>> & <<Nom>> & <<1 à 2 lignes>> \\
<<Repreneur 2>> & <<Nom>> & <<1 à 2 lignes>> \\
\bottomrule
\end{tabular}

\subsection{Tier 2 --- outreach J+3}

<<Idem, tableau de candidats secondaires.>>

\section{Données à exiger en data room}

\begin{enumerate}[leftmargin=2em]
  \item <<Élément 1>>
  \item <<Élément 2>>
  \item <<Élément 3>>
\end{enumerate}

\section{Contacts clés}

\textbf{Administrateur judiciaire.} <<Nom, cabinet, adresse, téléphone, e-mail.>>

\textbf{Mandataire judiciaire.} <<Idem.>>

\textbf{Conseil repreneur cité.} <<Nom, cabinet.>>

\textbf{Dirigeant historique.} <<Nom, fonction, LinkedIn.>>

\section*{Annexe --- Sources}
\addcontentsline{toc}{section}{Annexe --- Sources}

\begin{enumerate}[leftmargin=2em]
  \item <<Source 1 --- titre, URL, date de consultation.>>
  \item <<Source 2 --- idem.>>
\end{enumerate}

\textbf{Limitations méthodologiques.} <<Lister les sources non consultées,
les bases payantes en panne de crédits, les comptes non publiés au greffe,
toute hypothèse non vérifiable.>>

\end{document}
